% concatenated from v22fpt.tex
\documentclass{article}

\usepackage[utf8]{inputenc}
\usepackage{authblk}
\usepackage{hyperref}
\usepackage{enumitem}

\usepackage[a4paper, total={5in, 9in}]{geometry}

\usepackage{array}

%\setlength{\oddsidemargin}{-18mm}\setlength{\evensidemargin}{-18mm}

%\setlength{\oddsidemargin}{-5mm} %\setlength{\evensidemargin}{-5mm}   %%
\setlength{\topmargin}{-20mm}
%\setlength{\textwidth}{194mm} \setlength{\textheight}{270mm}

\usepackage{tikz}
\usepackage{tikz-cd}
\usetikzlibrary{arrows.meta}
\tikzcdset{arrow style=tikz,
diagrams={>={Stealth[round,length=4pt,width=6pt,inset=3pt]}}}

\usepackage{amsthm}

\usepackage{amsmath}
\usepackage{amssymb}
%\usepackage{bbold}
\usepackage{amsfonts}
\usepackage{faktor}
\usepackage[T1]{fontenc}
\usepackage{xfrac}
\usepackage{dsfont}
\usepackage{color}
\usepackage{todonotes}
\usepackage[english]{babel}
\usepackage{mathtools}

\newtheorem{theorem}{Theorem}
\newtheorem{definition}[theorem]{Definition}
\newtheorem{conjecture}[theorem]{Conjecture}
\newtheorem{lemma}[theorem]{Lemma}
\newtheorem{notation}[theorem]{Notation}
\newtheorem{observation}[theorem]{Observation}
\newtheorem{proposition}[theorem]{Proposition}
\newtheorem{property}[theorem]{Property}
\newtheorem{example}[theorem]{Example}
\newtheorem{remark}[theorem]{Remark}
\newtheorem{corollary}[theorem]{Corollary}

\newtheorem*{theorem*}{Theorem}

\DeclareMathOperator{\ord}{ord}
\DeclareMathOperator{\rank}{rank}
\DeclareMathOperator{\length}{length}
\DeclareMathOperator{\Hess}{Hess}
\DeclareMathOperator{\syz}{syz}
\DeclareMathOperator{\chara}{char}
\DeclareMathOperator{\jet}{jet}
\DeclareMathOperator{\Ann}{Ann}
\DeclareMathOperator{\Spec}{Spec}
\DeclareMathOperator{\Jac}{Jac}
\DeclareMathOperator{\fpt}{fpt}
\DeclareMathOperator{\lct}{lct}

\newcommand{\kx}{k[[\underline{x}]]}
\newcommand{\equivr}{\overset{\mathcal{R}}{\sim}}
\newcommand{\equivk}{\overset{\mathcal{K}}{\sim}}

\numberwithin{theorem}{section}


\usepackage{culmus}

\title{Bounds on the $F$-Pure Threshold of Isolated Hypersurface Singularities}
\author{Yotam Svoray}
\date{}

\AtEndDocument{\bigskip{\footnotesize%
  \textsc{Department of Mathematics, University of Utah, UT 84112} \par
  \textit{E-mail address}: \texttt{svoray@math.utah.edu}
}}

\begin{document}

\maketitle

\begin{abstract}
In this note, we obtain bounds for the $F$-pure threshold of isolated hypersurface singularities over an algebraically closed field of positive characteristic in terms of classical singularity invariants, notably the Milnor and Tjurina numbers. We also compare the $F$-pure threshold of a power series with that of its weighted initial form. For irreducible plane curves, this gives bounds in terms of the first two generators of the value semigroup, together with an equality criterion and an example showing that higher weighted-order terms can change the $F$-pure threshold. As applications, we derive bounds on the log canonical threshold and the Brian{\c{c}}on--Skoda exponent of complex isolated hypersurface singularities.
\end{abstract}

\section{Introduction}

The goal of this note is to use techniques from the study of isolated hypersurface singularities to bound and analyze the $F$-pure threshold (see Definition~\ref{def:Fpt}) of such singularities. Specifically, this note has three main results. First, we show how the Milnor and Tjurina numbers (see Definition~\ref{def:iso_sing}) of an isolated hypersurface singularity can be used to bound its $F$-pure threshold:

\begin{theorem}\label{ThmA}
Let $f\in\mathfrak{m}^2\subset k[[\underline{x}]]$ and assume that $\tau(f)<\infty$.
\begin{enumerate}
    \item\label{thm:fpt_milnor} If $\mu(f)<\infty$, then $\fpt(f)\geq\frac{1}{\tau(f)}-\frac{1}{\mu(f)}$.
    \item\label{cor:mil_infinite} If $\mu(f)=\infty$, then $\fpt(f)\geq\frac{1}{\tau(f)}$.
\end{enumerate}
\end{theorem}

Second, we show how weighted initial forms can be used to bound the $F$-pure threshold. In arbitrary dimension, if $g=\operatorname{in}_w(f)$ has weighted degree $d$ with respect to positive weights $w_1,\dots,w_n$, then $\fpt(g)\leq\fpt(f)\leq\min\{1,\frac{w_1+\cdots+w_n}{d}\}$. For an irreducible plane curve, this gives the following result in terms of the first two elements in the minimal generating sequence of its semigroup of values (see Definition~\ref{def:semigroup} and Lemma~\ref{lem:semi_prope}). The upper bound is inspired by the analogous formula for the log canonical threshold proved by Igusa in~\cite{igusa1977first} (see Remark~\ref{rem:igusa_lct}):

\begin{theorem}\label{thm:semigr_fpt}
Let $f\in\mathfrak{m}^2\subset k[[x,y]]$ be irreducible, and let $\beta_0,\beta_1$ be the first two elements of the minimal generating sequence of $\Gamma(f)$. Then:
\begin{enumerate}
    \item\label{cor:beta_bound} $\fpt(f)\leq\frac{1}{\beta_0}+\frac{1}{\beta_1}$.
    \item\label{cor:two_generators} If $\Gamma(f)=\langle\beta_0,\beta_1\rangle$, then
    \[
        \fpt(x^{\beta_0}+y^{\beta_1})\leq\fpt(f)\leq
        \frac{1}{\beta_0}+\frac{1}{\beta_1}.
    \]
    In particular, if $\fpt(x^{\beta_0}+y^{\beta_1})=\frac{1}{\beta_0}+\frac{1}{\beta_1}$, then $\fpt(f)=\fpt(x^{\beta_0}+y^{\beta_1})$.
\end{enumerate}
\end{theorem}

Third, we use these results to obtain analogous bounds for the log canonical threshold (see Definition~\ref{def:lct}) of a power series with integer coefficients. We also obtain a bound on the Brian{\c{c}}on--Skoda exponent (see Definition~\ref{def:eBS}) of such a power series:

\begin{corollary}\label{corB}
Let $f\in\mathfrak{m}^2\subset\mathbb{Z}[[\underline{x}]]$, viewed also as an element of $\mathbb{C}[[\underline{x}]]$, and assume that $\mu(f)<\infty$. Then:
\begin{enumerate}
    \item\label{cor:lct_milnor} $\lct(f)\geq\frac{1}{\tau(f)}-\frac{1}{\mu(f)}$.
    \item\label{cor:ebs} $\lct(f)\cdot\mu(f)+1\geq e^{\textup{BS}}(f)$.
\end{enumerate}
\end{corollary}

\noindent\textbf{Acknowledgments.} This work was done as part of the author's PhD thesis under the guidance of Karl Schwede, and we wish to thank him for his guidance, help, and support. We wish to thank Patricio Almir\'{o}n, Benjamin Baily, Gari Chua, Ilya Smirnov, and Jack Jeffries for their input on the ideas presented in this paper. In addition, we thank Gari Chua for his input and for reviewing several drafts of this note. The author was partially supported by NSF grant DMS-2101800.

\section{$F$-Pure Threshold of Isolated Singularities}\label{chap:F-stuff}

In this section, we show how the $F$-pure threshold of an isolated singularity can be related to its Milnor and Tjurina numbers. Several previous results compute and study the $F$-pure threshold of isolated singularities, especially for homogeneous and quasi-homogeneous polynomials with isolated singularities, as in~\cite{hernandez2016f,muller2018f,canton2016behavior,kadyrsizova2022lower}. For background on $F$-pure thresholds, we refer to~\cite{schwede2024singularities}.\\

Let $k$ be an algebraically closed field of characteristic $p>0$ and set $k[[\underline{x}]]=k[[x_1,\dots,x_n]]$, the ring of power series over $k$ in the variables $x_1,\dots,x_n$, with maximal ideal $\mathfrak{m}=\langle x_1,\dots,x_n\rangle$. We use multi-index notation when discussing power series.\\

We first recall the definitions and properties of isolated singularities used throughout this section:

\begin{definition}\label{def:iso_sing}
Let $f,g\in\mathfrak{m}\subset\kx$.
\begin{enumerate}
    \item \textbf{The Jacobian ideal} of $f$ is defined to be $\Jac(f)=\langle\partial_1(f),\dots,\partial_n(f)\rangle$, the ideal generated by the formal partial derivatives of $f$.
    \item The \textbf{Milnor} and \textbf{Tjurina} numbers of $f$ are, respectively,
    \[
        \mu(f)=\dim_k\!\left(\frac{\kx}{\Jac(f)}\right),\qquad
        \tau(f)=\dim_k\!\left(\frac{\kx}{\Jac(f)+\langle f\rangle}\right).
    \]
    \item We say that $f$ and $g$ are \textbf{right} (resp.\ \textbf{contact}) equivalent and denote $f\equivr g$ (resp.\ $f\equivk g$) if there exists a $k$-automorphism $\varphi$ of $\kx$ such that $\varphi(f)=g$ (resp.\ $\langle\varphi(f)\rangle=\langle g\rangle$).
    \item Given an ideal $I\subset\kx$ and $0\neq f\in\kx$, we denote by $\ord_I(f)$ the largest $N$ such that $f\in I^N$, and we write $\ord_{\mathfrak{m}}(f)=\ord(f)$. We set $\ord_I(0)=\infty$.
\end{enumerate}
\end{definition}

In the following proposition we summarize some known results about isolated singularities that we use throughout this text:

\begin{proposition}\label{prop:mt_basic}
Let $f\in\mathfrak{m}\subset\kx$. Then:
\begin{enumerate}
    \item If $\tau(f)<\infty$, then $\mu(f)<\infty$ if and only if $f\in\sqrt{\Jac(f)}$. In addition, if $f^N\in\Jac(f)$, then $\tau(f)\leq\mu(f)\leq N\cdot\tau(f)$.
    \item (Detecting Smoothness) $\mu(f)=0$ if and only if $\tau(f)=0$ if and only if $f\equivr x_1$.
    \item (Stability under $\equivk$) Given $g\in\kx$, if $f\equivk g$, then $\tau(f)=\tau(g)$.
    \item (Stability under $\equivr$) Given $g\in\kx$, if $f\equivr g$, then $\mu(f)=\mu(g)$.
    \item (Contact Determinacy) $\tau(f)<\infty$ if and only if there exists $N>0$ such that for every $g\in\kx$ with $g-f\in\mathfrak{m}^{N+1}$ we have $f\equivk g$. In this case we can choose $N=2\tau(f)-\ord(f)+2$.
    \item (Right Determinacy) $\mu(f)<\infty$ if and only if there exists $N>0$ such that for every $g\in\kx$ with $g-f\in\mathfrak{m}^{N+1}$ we have $f\equivr g$. In this case we can choose $N=2\mu(f)-\ord(f)+2$.
    \item If $\mu(f)<\infty$ (resp.\ $\tau(f)<\infty$), then $\mathfrak{m}^{\mu(f)}\subset\Jac(f)$ (resp.\ $\mathfrak{m}^{\tau(f)}\subset\Jac(f)+\langle f\rangle$).
\end{enumerate}
\end{proposition}

\begin{proof}
For the first item, see Lemma 4.21 in~\cite{svoray2025detecting}, which in turn generalizes results from~\cite{liu2018milnor}, and Proposition 2.1 in~\cite{hefez2019hypersurface}. For the second item, see Lemma 2.44 in Chapter I of~\cite{greuel2007introduction}. For the third and fourth items, see Lemma 1.2.7 in~\cite{boubakri2009hypersurface}. For the fifth and sixth items, see Corollary 2.4 and Theorem 2.8 in~\cite{boubakri2012invariants}; for the stated right-determinacy bound, note that their sharper bound $2\mu(f)-d(f)+1$, in terms of the differential order $d(f)$, is at most $2\mu(f)-\ord(f)+2$. For the seventh item, see Proposition 3.1.17 in~\cite{boubakri2009hypersurface}.
\end{proof}

\begin{remark}
\textup{With respect to the first item of Proposition~\ref{prop:mt_basic}, the quotient $\frac{\mu(f)}{\tau(f)}$ has been studied and bounded for singular power series with complex coefficients and finite Milnor number. For example, Liu in~\cite{liu2018milnor} showed that this quotient is bounded by $n$, the number of variables. In the case $n=2$, it was shown in~\cite{almiron2019note} that this quotient is bounded above by $\frac{4}{3}$, and for several families of surfaces with isolated singularities it was shown that this quotient is smaller than $\frac{3}{2}$. Finally, Ma and Zuo in~\cite{ma2026quotient} presented the optimal bound for this quotient in arbitrary dimension, with an analogous bound for a similar quotient in positive characteristic. For more information, see~\cite{almiron2022quotient,almiron20214,alberich2021minimal}. In addition, Saito in~\cite{saito1971quasihomogene} showed that for complex power series $\mu(f)-\tau(f)$ equals the $(n-1)$st Poincar\'e number, and therefore $\mu(f)=\tau(f)$ if and only if $f$ is right equivalent to a quasi-homogeneous polynomial.}
\end{remark}

\begin{lemma}\label{lem:jac_frob}
If $f\in\mathfrak{m}^{[p^e]}$, then $\Jac(f)\subset\mathfrak{m}^{[p^e]}$.
\end{lemma}

\begin{proof}
If $f\in\mathfrak{m}^{[p^e]}$, then there exist $a_1,\dots,a_n\in\kx$ such that $f=\sum_{j=1}^n a_jx_j^{p^e}$. Therefore, for every $i$ we have $\partial_i(f)=\sum_{j=1}^n(\partial_i(a_j)x_j^{p^e}+a_j\partial_i(x_j^{p^e}))$. Since $k$ has characteristic $p$, $\partial_i(x_j^{p^e})=0$ for every $i,j$. Thus $\partial_i(f)=\sum_{j=1}^n\partial_i(a_j)x_j^{p^e}\in\mathfrak{m}^{[p^e]}$, and the result follows.
\end{proof}

We now define the $F$-pure threshold of a power series:

\begin{definition}\label{def:Fpt}
Given a proper ideal $J\subsetneq\kx$ and $f\in\sqrt J$, we denote by $\nu_e(f,J)$ the largest integer $N\geq0$ such that $f^N\notin J^{[p^e]}$, and we set $c(f,J)=\lim_{e\to\infty}\frac{\nu_e(f,J)}{p^e}$. If $J=\mathfrak{m}$, we write $\nu_e(f)=\nu_e(f,\mathfrak{m})$ and $\fpt(f)=c(f,\mathfrak{m})$, which we call the \textbf{$F$-pure threshold} of $f$.
\end{definition}

\begin{remark}\label{rem:fpt_contact}
\textup{Observe that $\fpt(f)$ is stable under contact, and therefore also under right, equivalence. Indeed, a $k$-automorphism $\varphi$ of $\kx$ is local, so $\varphi(\mathfrak{m}^{[p^e]})=\mathfrak{m}^{[p^e]}$, and multiplication by a unit does not affect membership in $\mathfrak{m}^{[p^e]}$.}
\end{remark}

\begin{example}
\textup{Assume that $p>2$ and let $f\in\kx$ be of order $2$. By the splitting lemma (see, for example, Theorem 2.1 in~\cite{greuel2025splitting}), we have $f\equivr x_1^2+\cdots+x_r^2+g(x_{r+1},\dots,x_n)$. If $r\geq2$, then by Fedder's criterion (see~\cite{fedder1983F} or Chapter 4.1 in~\cite{schwede2024singularities}) we have $\nu_e(f)=p^e-1$, and hence $\fpt(f)=1$.}
\end{example}

The following lemma is based on the proof of Lemma 3.3 in~\cite{schwede2024singularities} and the proof of Proposition 1.9 in~\cite{mustata2005f}:

\begin{lemma}\label{lem:bound_c_nu}
Given a proper ideal $I\subsetneq\kx$ and $f\in\sqrt I$, for every $e$ we have
\[
    p\nu_e(f,I)\leq\nu_{e+1}(f,I)\leq p(\nu_e(f,I)+1)-1.
\]
In particular, $\nu_0(f,I)\leq c(f,I)\leq\nu_0(f,I)+1$.
\end{lemma}

\begin{proof}
Since $\kx$ is regular, Frobenius is faithfully flat. Hence $f^{\nu_e(f,I)}\notin I^{[p^e]}$ implies $f^{p\nu_e(f,I)}\notin I^{[p^{e+1}]}$. Therefore, $p\nu_e(f,I)\leq\nu_{e+1}(f,I)$, and by induction $p^e\nu_0(f,I)\leq\nu_e(f,I)$. Similarly, since $f^{\nu_e(f,I)+1}\in I^{[p^e]}$, we have $f^{p(\nu_e(f,I)+1)}\in I^{[p^{e+1}]}$. Hence $p(\nu_e(f,I)+1)\geq\nu_{e+1}(f,I)+1$. Thus $\nu_0(f,I)\leq\frac{\nu_e(f,I)}{p^e}\leq\nu_0(f,I)+1$, and the result follows by taking $e\to\infty$.
\end{proof}

\begin{example}
\textup{Assume $p>5$ and let $f\in\mathfrak{m}^2\subset k[[x,y]]$. By~\cite{greuel1990simple}, either $f\in\langle x,y^2\rangle^3$, up to contact equivalence, or $f$ is contact equivalent to an ADE singularity. In the first case, assigning weights $2$ and $1$ to $x$ and $y$, respectively, gives $\fpt(f)\leq\frac12$ by the weighted-degree estimate of Proposition~\ref{prop:weighted_fpt}. On the other hand, direct computation for ADE singularities (see Example 2.5 in~\cite{takagi2004f}) gives $\fpt(f)>\frac12$. Hence $f$ is contact equivalent to an ADE singularity if and only if $\fpt(f)>\frac12$.}
\end{example}

\begin{proposition}\label{prop:det_fpt_m-adic}
Let $0\neq f\in\mathfrak{m}^2$. Then there exists $N\gg0$ such that for every $g\in(\Jac(f)+\langle f\rangle)^N$ we have $\fpt(f)=\fpt(f+g)$.
\end{proposition}

\begin{proof}
By Theorem 2.4 and Corollary 2.22 in~\cite{svoray2025detecting}, there exists $N$ such that for every $g\in(\Jac(f)+\langle f\rangle)^N$ we have $f\equivk f+g$. Therefore, by Remark~\ref{rem:fpt_contact}, $\fpt(f)=\fpt(f+g)$.
\end{proof}

\begin{remark}
\begin{enumerate}
    \item \textup{Proposition~\ref{prop:det_fpt_m-adic} is a generalization of Theorem 4.3 in~\cite{hernandez2018local}, where it is assumed that $f$ has an isolated singularity and hence $\sqrt{\langle f\rangle+\Jac(f)}=\mathfrak{m}$. In addition, by Proposition~\ref{prop:mt_basic}, if $\tau(f)<\infty$, then contact determinacy gives the explicit safe choice $N=2\tau(f)-\ord(f)+3$: indeed, $\Jac(f)+\langle f\rangle\subset\mathfrak{m}$, so its $N$th power is contained in $\mathfrak{m}^{N}$, which is deep enough for $(N-1)$-determinacy. This still improves the bound given in~\cite{hernandez2018local}, namely $N=p^{2\tau(f)}\cdot n$.\footnote{We note that the author of~\cite{hernandez2018local}, Hern\'{a}ndez, was informed of the potential application of determinacy to $\mathfrak{m}$-adic constancy of $\fpt$ by Holger Brenner in a private conversation.}}
    \item \textup{The version of Proposition~\ref{prop:det_fpt_m-adic} obtained by replacing $\Jac(f)+\langle f\rangle$ with $\mathfrak{m}$ is not true if $f$ does not define an isolated singularity. For example, over $k[[x,y]]$, if $d$ is coprime to $p$, then by Theorem 3.1 of~\cite{hernandez2015F}, $\fpt(x^d+y^n)>\fpt(x^d)=\frac1d$ for every $n\gg d$. What we do know is that $\fpt$ is lower semicontinuous in the $\mathfrak{m}$-adic topology: for every $f\in\mathfrak{m}$ there exists $N$ such that for every $g\in\mathfrak{m}^N$ we have $\fpt(f+g)\geq\fpt(f)$. This follows from the ascending chain conjecture, first stated as Conjecture 4.4 in~\cite{blickle2009F} and proved by Sato in~\cite{sato2019ascending}.}
\end{enumerate}
\end{remark}

\begin{lemma}\label{lem:c(f,Jac)_fpt}
Let $0\neq f\in\mathfrak{m}$. Then:
\begin{enumerate}
    \item If $\mu(f)<\infty$, then $\fpt(f)\mu(f)\geq c(f,\Jac(f))$.
    \item If $0<\tau(f)<\infty$, then $\fpt(f)\tau(f)\geq c(f,\Jac(f)+\langle f\rangle)$.
\end{enumerate}
\end{lemma}

\begin{proof}
If $\mu(f)<\infty$, Proposition~\ref{prop:mt_basic} gives $\mathfrak{m}^{\mu(f)}\subset\Jac(f)$, and hence
\[
(\mathfrak{m}^{[p^e]})^{\mu(f)}\subset\Jac(f)^{[p^e]}.
\]
Since $f^{\nu_e(f)+1}\in\mathfrak{m}^{[p^e]}$, it follows that
\[
\mu(f)(\nu_e(f)+1)\geq\nu_e(f,\Jac(f))+1.
\]
Dividing by $p^e$ and taking the limit proves the first assertion.  Similarly, if $0<\tau(f)<\infty$, then $\mathfrak{m}^{\tau(f)}\subset\Jac(f)+\langle f\rangle$, so
\[
(\mathfrak{m}^{[p^e]})^{\tau(f)}
\subset(\Jac(f)+\langle f\rangle)^{[p^e]}.
\]
The same argument gives
\[
\tau(f)(\nu_e(f)+1)
\geq\nu_e(f,\Jac(f)+\langle f\rangle)+1,
\]
and the second assertion follows after dividing by $p^e$ and taking the limit.
\end{proof}

\begin{corollary}\label{cor:BS_F}
If $0<\mu(f)<\infty$, then $\fpt(f)\geq\frac{\nu_0(f,\Jac(f))}{\mu(f)}$.
\end{corollary}

\begin{proof}
This follows from Lemmas~\ref{lem:bound_c_nu} and~\ref{lem:c(f,Jac)_fpt}.
\end{proof}

\begin{corollary}
If $0<\tau(f)<\infty$, then $\fpt(f)\geq\frac{1}{\tau(f)}\left(1-\frac{1}{\ord(f)}\right)$.
\end{corollary}

\begin{proof}
Note that $f\notin\mathfrak{m}^{\ord(f)+1}$ and $\Jac(f)\subset\mathfrak{m}^{\ord(f)-1}$. Therefore, from item 3 of Proposition 1.7 in~\cite{mustata2005f} $c(f,\Jac(f)+\langle f\rangle)$ is bounded below by $1-\frac{1}{\ord(f)}$, and the result follows from Lemma~\ref{lem:c(f,Jac)_fpt}.
\end{proof}

\begin{proof}[Proof of Part~\ref{thm:fpt_milnor} of Theorem~\ref{ThmA}]
This follows from Lemmas~\ref{lem:c(f,Jac)_fpt} and~\ref{lem:bound_c_nu}, together with item 1 of Proposition~\ref{prop:mt_basic}, which gives $\frac{\mu(f)}{\tau(f)}\leq\nu_0(f,\Jac(f))+1$.
\end{proof}

Part~\ref{thm:fpt_milnor} of Theorem~\ref{ThmA} leads directly to the question of what happens if $\tau(f)$ is finite but $\mu(f)$ is not. In this case, regarding $\frac{1}{\mu(f)}$ as $0$, one would expect $\fpt(f)\geq\frac{1}{\tau(f)}$. This is, in fact, true and appears as Part~\ref{cor:mil_infinite} of Theorem~\ref{ThmA}. It follows from the following lemma, whose proof is inspired by the proof of Proposition 3.3 in~\cite{pellikaan1989series}:

\begin{lemma}\label{lem:milnor_limit}
Let $f\in\kx$ such that $\mu(f)=\infty$ and $\tau(f)<\infty$. Then there exist $a\in\mathfrak{m}\setminus\mathfrak{m}^2$ and $l_0\gg0$ such that for every $l>l_0$ not divisible by $p$ we have $\mu(f+a^l)<\infty$.
\end{lemma}

\begin{proof}
Since $\tau(f)<\infty$, the ideal $\langle f\rangle+\Jac(f)$ is $\mathfrak{m}$-primary, while $\mu(f)=\infty$ implies that $\frac{\kx}{\Jac(f)}$ is one-dimensional. By Theorem 1.7 in~\cite{ulrich1987theory}, alternatively by~\cite{huneke1987structure}, after a linear change of coordinates we may assume that $\mathfrak{a}=\langle\partial_2(f),\dots,\partial_n(f)\rangle$ has height $n-1$. Hence $\frac{\kx}{\mathfrak{a}}$ is a one-dimensional complete intersection. We may also choose the coordinates so that $\mathfrak{a}+\langle x_1\rangle$ is $\mathfrak{m}$-primary. Let $\mathfrak{p}_1,\dots,\mathfrak{p}_r$ be the associated primes of $\frac{\kx}{\mathfrak{a}}$; these are precisely its minimal primes. For a fixed $\mathfrak{p}_i$ there is at most one $l$, with $p\nmid l$, such that $\partial_1(f)+lx_1^{l-1}\in\mathfrak{p}_i$. Indeed, if this holds for $l_1>l_2$, then
\[
x_1^{l_2-1}(l_1x_1^{l_1-l_2}-l_2)\in\mathfrak{p}_i.
\]
Since $x_1\notin\mathfrak{p}_i$ and $l_2$ is nonzero in $k$, the second factor is a unit, a contradiction. Thus, for every sufficiently large $l$ not divisible by $p$, the element $\partial_1(f)+lx_1^{l-1}$ is a nonzerodivisor on $\frac{\kx}{\mathfrak{a}}$. Consequently $\Jac(f+x_1^l)=\mathfrak{a}+\langle\partial_1(f)+lx_1^{l-1}\rangle$ is $\mathfrak{m}$-primary, and hence $\mu(f+x_1^l)<\infty$.
\end{proof}

\begin{remark}\label{rem:limit_milnor}
\textup{In Lemma~\ref{lem:milnor_limit}, we must have $\mu(f+a^l)\to\infty$ as $l\to\infty$, where $l$ is not divisible by $p$. Otherwise the set $\{2\mu(f+a^l)-\ord(f)+2\}_{l>l_0}$ would be bounded; let $N$ be an upper bound. Since $\ord(f)=\ord(f+a^l)$ for $l\gg0$, right determinacy implies that each $f+a^l$ is $N$-determined. Since $f+a^l\to f$ in the $\mathfrak{m}$-adic topology, for $l\gg0$ we obtain $f+a^l\equivr f$, and hence $\mu(f)=\mu(f+a^l)<\infty$, a contradiction.}
\end{remark}

\begin{proof}[Proof of Part~\ref{cor:mil_infinite} of Theorem~\ref{ThmA}]
By Lemma~\ref{lem:milnor_limit}, there exists $a$ such that $f_l=f+a^l$ satisfies $\mu(f_l)<\infty$ for all sufficiently large $l$ not divisible by $p$. By contact determinacy, $f_l\equivk f$ for all sufficiently large such $l$, because $a^l$ has arbitrarily high order. Hence $\tau(f_l)=\tau(f)$ and, by Remark~\ref{rem:fpt_contact}, $\fpt(f_l)=\fpt(f)$. Part~\ref{thm:fpt_milnor} of Theorem~\ref{ThmA} therefore gives
\[
    \fpt(f)\geq\frac{1}{\tau(f)}-\frac{1}{\mu(f_l)}.
\]
By Remark~\ref{rem:limit_milnor}, $\mu(f_l)\to\infty$ as $l\to\infty$ through integers not divisible by $p$. Taking the limit yields $\fpt(f)\geq\frac{1}{\tau(f)}$.
\end{proof}

\section{$F$-Pure Threshold and Semigroup of Values}

In this section, we show how the generators of the semigroup of values of an irreducible power series in two variables over $k$ can be used to bound its $F$-pure threshold. Throughout this section, we assume that $n=2$ and denote $k[[x_1,x_2]]$ by $k[[x,y]]$.

\begin{definition}\label{def:semigroup}
Given an irreducible $f\in k[[x,y]]$, the \textbf{semigroup of values} of $f$, denoted $\Gamma(f)$, is the set of all values
\[
i(f,g)=\dim_k\left(\frac{k[[x,y]]}{\langle f,g\rangle}\right),
\]
where $g\in k[[x,y]]\setminus\langle f\rangle$.
\end{definition}

We summarize the main results about $\Gamma(f)$ that we use throughout this section in the following lemma:

\begin{lemma}\label{lem:semi_prope}
Let $f\in k[[x,y]]$ be an irreducible power series. Then:
\begin{enumerate}
    \item If $g,h\notin\langle f\rangle$, then $i(f,gh)=i(f,g)+i(f,h)$.
    \item If $g\notin\langle f\rangle$, then $i(f,g)\geq\ord(f)\ord(g)$.
    \item $\Gamma(f)$ is a numerical semigroup, i.e.\ $\gcd(\Gamma(f))=1$.
    \item $\min(\Gamma(f)\setminus\{0\})=\ord(f)$.
    \item There exists a finite sequence of integers $\beta_0,\dots,\beta_K$ that generates $\Gamma(f)$ as a semigroup and is minimal under inclusion.
    \item $\Gamma(f)$ is generated by two integers $\alpha$ and $\beta$ if and only if $f$ is contact equivalent to a power series of the form
    \[
        x^\alpha+y^\beta+
        \sum_{\alpha j+\beta i>\alpha\beta}a_{i,j}x^iy^j
    \]
    for some $a_{i,j}\in k$.
    \item $f$ is contact equivalent to $y^{\beta_0}+\sum_{l=0}^{\beta_0-1}c_l(x)y^l$, where $c_l(x)\in k[[x]]$ satisfies $\ord(c_l(x))\geq\frac{(\beta_0-l)\beta_1}{\beta_0}$ for every $l>0$ and $\ord(c_0(x))=\beta_1$.
\end{enumerate}
\end{lemma}

\begin{proof}
For item 1, see, for example, item 6 of Lemma 3.3 in~\cite{svoray2026Milnor}; for item 2, see Theorem 3 of Chapter 3.3 in~\cite{fulton2008algebraic}; for item 3, see Lemma 3.1 in~\cite{barroso2012approach}; and for item 4, see item 5 of Lemma 3.3 in~\cite{svoray2026Milnor}. For item 5, set $\beta_0(f)=\min(\Gamma(f)\setminus\{0\})$ and $\beta_{i+1}(f)=\min(\Gamma(f)\setminus(\beta_0\mathbb{N}+\cdots+\beta_i\mathbb{N}))$. By definition, the sequence $\beta_0,\beta_1,\dots$ generates $\Gamma(f)$ and is minimal under inclusion, while Proposition 6.1 in Chapter 6 of~\cite{hefez2003irreducible} shows that this sequence terminates. For item 6 in arbitrary characteristic, see the description of branches with a prescribed two-generated semigroup in Section 10, in particular Corollary 10.7, of~\cite{barroso2012approach}; compare also Proposition 2.1 in Chapter 6 of~\cite{zariski2006moduli}. Finally, item 7 follows from the Weierstrass  preparation theorem (see, for example,~\cite{o1972weierstrass}) and the Semigroup Theorem (see Theorem 3.2 in~\cite{barroso2012approach}).
\end{proof}

\begin{remark}\label{rem:2dim_sing}
\begin{enumerate}
    \item \textup{If $f\in k[[x,y]]$ is irreducible, then $\tau(f)<\infty$, since $f$ defines a reduced plane curve whose singular locus is supported at the closed point. This need not be true for the Milnor number. For example, if $\operatorname{char}(k)=p$, then $\mu(x^p+y^{p+1})=\infty$ but $\tau(x^p+y^{p+1})=p^2$.}
    \item \textup{By item 7 of Lemma~\ref{lem:semi_prope}, up to contact equivalence we have $i(f,x)=\beta_0=\ord(f)$. This gives an alternative proof of $\fpt(f)\geq \frac{1}{\ord(f)}$ (see Proposition 4.1 in~\cite{takagi2004f}). Indeed, from $f^{\nu_e(f)+1}\in\langle x,y\rangle^{[p^e]}$ we may write $f^{\nu_e(f)+1}=gx^{p^e}+hy^{p^e}$. Reducing modulo $x$ gives $(\nu_e(f)+1)\ord(f)\geq p^e$, and hence $\frac{1}{\ord(f)}\leq\frac{\nu_e(f)+1}{p^e}\to\fpt(f)$.}
    \item \textup{The semigroup of values is an important tool in the study of curve singularities and of one-dimensional local rings more generally. For more information, see, for example,~\cite{assi2020numerical,barroso2012approach,barroso2022note,castellanos2005semigroup,campillo1994gorenstein,d2025value,de1987semigroup,kunz1970value,zariski2006moduli}.}
    \item \textup{In item 5 of Lemma~\ref{lem:semi_prope}, $K>0$ if and only if $f\in\mathfrak{m}^2$, since a numerical semigroup has a unique generator if and only if it is $\mathbb{N}$, which is equivalent to $\ord(f)=1$.}
\end{enumerate}
\end{remark}

\begin{definition}\label{def:weighted_initial}
Let $w=(w_1,\dots,w_n)\in\mathbb{N}_{>0}^n$. For $\underline{x}^a=x_1^{a_1}\cdots x_n^{a_n}$, define $\operatorname{wt}_w(\underline{x}^a)=\sum_iw_i a_i$. If $0\neq f=\sum_a c_a\underline{x}^a\in k[[\underline{x}]]$, let $\ord_w(f)=\min\{\operatorname{wt}_w(\underline{x}^a):c_a\neq0\}$ and let $\operatorname{in}_w(f)$ be the sum of the terms of $f$ having $w$-weight $\ord_w(f)$.
\end{definition}

\begin{proposition}\label{prop:weighted_fpt}
Let $0\neq f\in\mathfrak{m}\subset k[[x_1,\dots,x_n]]$, let $w=(w_1,\dots,w_n)\in\mathbb{N}_{>0}^n$, set $d=\ord_w(f)$ and $g=\operatorname{in}_w(f)$. Then
\[
    \fpt(g)\leq\fpt(f)\leq
    \min\left\{1,\frac{w_1+\cdots+w_n}{d}\right\}.
\]
\end{proposition}

\begin{proof}
Write $f=g+h$, where every monomial of $h$ has $w$-weight greater than $d$, and fix $q=p^e$. If $g^N\notin\mathfrak{m}^{[q]}$, then some monomial $\underline{x}^a$ of $g^N$ has nonzero coefficient and $a_i<q$ for every $i$. It has $w$-weight $Nd$, while every monomial of $f^N-g^N$ has $w$-weight greater than $Nd$; hence it cannot cancel in $f^N$. Thus $\nu_e(f)\geq\nu_e(g)$ for every $e$, and $\fpt(f)\geq\fpt(g)$. Conversely, every monomial of $f^N$ has $w$-weight at least $Nd$. A monomial outside $\mathfrak{m}^{[q]}$ has exponents at most $q-1$, and therefore has $w$-weight at most $(q-1)(w_1+\cdots+w_n)$. Hence $Nd>(q-1)(w_1+\cdots+w_n)$ implies $f^N\in\mathfrak{m}^{[q]}$. Dividing the resulting bound on $\nu_e(f)$ by $q$ and letting $e\to\infty$ gives $\fpt(f)\leq \frac{w_1+\cdots+w_n}{d}$. Finally, $\fpt(f)\leq1$ since $f^q\in\mathfrak{m}^{[q]}$.
\end{proof}

\begin{remark}\label{rem:weighted_equality}
\textup{In particular, if $\fpt(\operatorname{in}_w(f))=\frac{w_1+\cdots+w_n}{\ord_w(f)}\leq1$, then Proposition~\ref{prop:weighted_fpt} gives $\fpt(f)=\fpt(\operatorname{in}_w(f))$.}
\end{remark}

\begin{proof}[Proof of Part~\ref{cor:beta_bound} of Theorem~\ref{thm:semigr_fpt}]
By item 7 of Lemma~\ref{lem:semi_prope}, up to contact equivalence we may write $f=y^{\beta_0}+\sum_{l=0}^{\beta_0-1}c_l(x)y^l$, where $\ord(c_l(x))\geq\frac{(\beta_0-l)\beta_1}{\beta_0}$ and $\ord(c_0(x))=\beta_1$. Give $x$ weight $\beta_0$ and $y$ weight $\beta_1$. Then $\ord_w(f)=\beta_0\beta_1$, so Proposition~\ref{prop:weighted_fpt} gives $\fpt(f)\leq \frac{\beta_0+\beta_1}{\beta_0\beta_1}=\frac{1}{\beta_0}+\frac{1}{\beta_1}$.
\end{proof}

\begin{remark}\label{rem:igusa_lct}
\textup{Igusa in Theorem 1 of~\cite{igusa1977first} proved that if $f\in\mathbb{C}[[x,y]]$ is irreducible, then $\lct(f)=\frac{1}{\beta_0}+\frac{1}{\beta_1}$. An alternative proof was provided by Kuwata in Theorem 1.1 of~\cite{kuwata1999log}. This result was later generalized by Lee in~\cite{lee2026log} to arbitrary algebraically closed fields. Thus Part~\ref{cor:beta_bound} of Theorem~\ref{thm:semigr_fpt} can be viewed as an $F$-pure threshold analogue of Igusa's result.}
\end{remark}

\begin{proof}[Proof of Part~\ref{cor:two_generators} of Theorem~\ref{thm:semigr_fpt}]
By item 6 of Lemma~\ref{lem:semi_prope}, up to contact equivalence we may write $f=x^{\beta_0}+y^{\beta_1}+h$, where every monomial $x^iy^j$ of $h$ satisfies $\beta_0j+\beta_1i>\beta_0\beta_1$. Giving $x$ weight $\beta_1$ and $y$ weight $\beta_0$, we have $\operatorname{in}_w(f)=x^{\beta_0}+y^{\beta_1}$ and $\ord_w(f)=\beta_0\beta_1$. Proposition~\ref{prop:weighted_fpt} gives
\[
    \fpt(x^{\beta_0}+y^{\beta_1})\leq\fpt(f)
    \leq\frac1{\beta_0}+\frac1{\beta_1}.
\]
The final assertion follows when the two endpoints agree.
\end{proof}

\begin{example}\label{ex:higher_terms_change_fpt}
\textup{The equality $\fpt(f)=\fpt(x^{\beta_0}+y^{\beta_1})$ need not hold in general. Let $k=\overline{\mathbb{F}}_2$ and $f=x^5+y^7+x^3y^3$. Give $x,y$ weights $7,5$, respectively. Then $\operatorname{in}_w(f)=x^5+y^7$, which is irreducible since $\gcd(5,7)=1$. Any nontrivial factorization of $f$ would induce a nontrivial factorization of its weighted initial form, so $f$ is irreducible. By item 6 of Lemma~\ref{lem:semi_prope}, $\Gamma(f)=\langle5,7\rangle$. By Theorem 3.1 of~\cite{hernandez2015F}, we have $\fpt(x^5+y^7)=\frac{1}{4}$. Put $I_q=\langle x^q,y^q\rangle$, where $q=2^e$. If $q=4^m$, induction gives}
\[
f^{\frac{q-1}{3}}\equiv x^{q-1}y^{q-1}\pmod{I_q}.
\]
\textup{Indeed, if the congruence holds for $q$, then, in characteristic $2$,}
\[
f^{\frac{4q-1}{3}}
=f\big(f^{\frac{q-1}{3}}\big)^4
\equiv(x^5+y^7+x^3y^3)x^{4q-4}y^{4q-4}
\equiv x^{4q-1}y^{4q-1}\pmod{I_{4q}}.
\]
\textup{If instead $q=2\cdot4^m$, set $r=\frac{q}{2}=4^m$. Squaring the congruence for $r$ gives}
\[
f^{\frac{q-2}{3}}\equiv x^{q-2}y^{q-2}\pmod{I_q}.
\]
\textup{Thus, in either case, $f^{\lfloor\frac{q-1}{3}\rfloor}\notin I_q$, while multiplying the displayed surviving monomial by each term of $f$ shows that $f^{\lfloor\frac{q-1}{3}\rfloor+1}\in I_q$. Hence $\nu_e(f)=\lfloor\frac{q-1}{3}\rfloor$ for every $e$, and therefore $\fpt(f)=\frac13>\frac14=\fpt(x^5+y^7)$.}
\end{example}

\begin{remark}
\begin{enumerate}
    \item \textup{If $p\nmid\alpha\beta$, Theorem 3.1 of~\cite{hernandez2015F} computes $\fpt(x^\alpha+y^\beta)$ from the base-$p$ expansions of $\frac{1}{\alpha}$ and $\frac{1}{\beta}$. Whenever $\fpt(x^\alpha+y^\beta)=\frac{1}{\alpha}+\frac{1}{\beta}$, Remark~\ref{rem:weighted_equality} shows that every perturbation by terms of strictly larger $(\beta,\alpha)$-weight has the same $F$-pure threshold.}
    \item \textup{Arrigoni in~\cite{arrigoni2024f} studied the $F$-pure threshold of $f\in k[[x,y]]$, and more generally $c(f,\langle x^a,y^b\rangle)$, when $\Gamma(f)$ is generated by two elements, calling these ``simple algebroid plane branches.''}
\end{enumerate}
\end{remark}

\section{Log Canonical Threshold and Reduction mod $p$}

In this section we show how the previous results can be applied to the log canonical threshold and conclude Corollary~\ref{corB}. For more information on log canonical thresholds, see Chapter 9.3B in~\cite{lazarsfeld2003positivity} or~\cite{kollar2008powers,blum2016log}. The connection between contact equivalence, the Milnor number of a complex power series, and the log canonical threshold has also been studied, for example, in~\cite{almiron2022limit,cluckers2022log,mustactua2026some}.

\begin{definition}\label{def:lct}
Given a nonzero power series $f\in\mathfrak{m}\subset\mathbb{C}[[\underline{x}]]$, the \textbf{log canonical threshold} of $f$, denoted $\lct(f)$, is
\[
    \lct(f)=\inf_E\frac{A(E)}{\ord_E(f)},
\]
where $E$ runs over divisorial valuations centered at the origin and $A(E)$ denotes the log discrepancy of $E$.
\end{definition}

\begin{remark}
\textup{If $f$ is convergent, this agrees with the analytic characterization
\[
\lct(f)=\sup\left\{\lambda\geq0:\ |f|^{-\lambda}\text{ is locally }L^2
\text{ near }0\in\mathbb{C}^n\right\}.
\]
Equivalently, $\int_U |f|^{-2\lambda}\,d\mu<\infty$ for some neighborhood $U$ of the origin.}
\end{remark}


For the proof of Corollary~\ref{corB}, right determinacy allows us to replace a power series $f\in\mathbb{Z}[[\underline{x}]]$ with $\mu(f)<\infty$ by a sufficiently high polynomial truncation, without changing the invariants appearing in the statement. Thus, from now until the end of this section, let $f\in\mathbb{Z}[\underline{x}]$ be a polynomial with integer coefficients, and for every prime $p$ let $f_p\in\overline{\mathbb{F}}_p[[\underline{x}]]$ denote its reduction modulo $p$, followed by extension of scalars to $\overline{\mathbb{F}}_p$. The following proposition, presented in~\cite{hara2003generalization}, shows the relation between $\lct(f)$ and $\fpt(f_p)$. For more information on this relation, see Chapter 6.4 in~\cite{schwede2024singularities} or~\cite{mustata2012impanga}.

\begin{proposition}\label{prop:lct_fpt}
For every sufficiently large prime $p$, we have $\fpt(f_p)\leq\lct(f)$. Moreover, as $p$ ranges over the primes,
\[
    \lim_{p\to\infty}\fpt(f_p)=\lct(f).
\]
\end{proposition}

An analogous result for the Milnor and Tjurina numbers is presented as Corollary 55 in~\cite{greuel2021semicontinuity}. Given $f\in\mathbb{Z}[[\underline{x}]]$, we denote by $\mu_p(f)$ and $\tau_p(f)$ the Milnor and Tjurina numbers of $f_p$, and by $\mu(f)$ and $\tau(f)$ the Milnor and Tjurina numbers of $f\in\mathbb{C}[[\underline{x}]]$, respectively.

\begin{proposition}\label{prop:tau_p}
Let $f\in\mathbb{Z}[[\underline{x}]]$. Then:
\begin{enumerate}
    \item If $\mu(f)<\infty$, then $\mu(f)=\mu_p(f)$ for every $p\gg0$.
    \item If $\mu_p(f)<\infty$ for some $p$, then $\mu(f)\leq\mu_p(f)$ and $\mu_q(f)\leq\mu_p(f)$ for every $q\gg0$.
\end{enumerate}
\end{proposition}

\begin{remark}\label{rem:tau_p}
\textup{An analogous result to Proposition~\ref{prop:tau_p} holds for $\tau(f)$ and $\tau_p(f)$, with the proof appearing in the same reference.}
\end{remark}

\begin{proof}[Proof of Part~\ref{cor:lct_milnor} of Corollary~\ref{corB}]
Since $\mu(f)<\infty$, we also have $\tau(f)<\infty$. Proposition~\ref{prop:tau_p} and Remark~\ref{rem:tau_p} imply that, for every $p\gg0$, $\mu_p(f)=\mu(f)$ and $\tau_p(f)=\tau(f)$. Therefore, Theorem~\ref{ThmA} gives
\[
    \lct(f)\geq\fpt(f_p)\geq
    \frac{1}{\tau_p(f)}-\frac{1}{\mu_p(f)}
    =\frac{1}{\tau(f)}-\frac{1}{\mu(f)},
\]
and the result follows.
\end{proof}

Similarly, we present an analogous result for Corollary~\ref{cor:BS_F}.

\begin{definition}\label{def:eBS}
Given $f\in\mathfrak{m}\subset\mathbb{C}[[\underline{x}]]$ such that $f\in\sqrt{\Jac(f)}$, the \textbf{Brian{\c{c}}on--Skoda exponent} of $f$, denoted $e^{\textup{BS}}(f)$, is the smallest integer $N\geq1$ such that $f^N\in\Jac(f)$.
\end{definition}

\begin{remark}
\textup{The name Brian{\c{c}}on--Skoda comes from the Brian{\c{c}}on--Skoda theorem (see~\cite{skoda1974cloture}), which gives $e^{\textup{BS}}(f)\leq n$ for every $f\in\mathfrak{m}\subset\mathbb{C}[[\underline{x}]]$, where $n$ is the number of variables. Similarly, by the Brian{\c{c}}on--Skoda theorem for power series in positive characteristic, as explained in Remark 4.5 of~\cite{hefez2019hypersurface}, if $f\in\sqrt{\Jac(f)}$, then $\nu_0(f,\Jac(f))+1\leq n$. The Brian{\c{c}}on--Skoda exponent has been studied in the context of topology and complex isolated singularities, for example, in~\cite{jung2022brianccon} and in Section 3 of~\cite{almiron2023tjurina}, which in turn is based on~\cite{varvcenko1982asymptotic}.}
\end{remark}

\begin{lemma}\label{lem:nu_E_BS}
Assume $\mu(f)<\infty$. Then, for every sufficiently large prime $p$,
\[
\nu_0(f_p,\Jac(f_p))+1=e^{\textup{BS}}(f).
\]
\end{lemma}

\begin{proof}
Set $N=e^{\textup{BS}}(f)$. Because extension of scalars from $\mathbb{Q}$ to $\mathbb{C}$ is faithfully flat, both the Milnor number and the membership conditions $f^j\in\Jac(f)$ may be computed over $\mathbb{Q}$. Since $\mu(f)<\infty$, the Jacobian algebra $A=\frac{\mathbb{Q}[[\underline{x}]]}{\Jac(f)}$ has finite length. After truncating modulo a sufficiently high power of the maximal ideal, $A$ and the class of $f$ are defined over $\mathbb{Z}[\frac{1}{c}]$ for some $c\neq0$. After enlarging $c$ and applying generic freeness, the relations $f^N=0$ and $f^{N-1}\neq0$ in $A$ remain true after reduction modulo every sufficiently large prime. Hence $f_p^N\in\Jac(f_p)$ and $f_p^{N-1}\notin\Jac(f_p)$ for every $p\gg0$, which is exactly $\nu_0(f_p,\Jac(f_p))+1=N$.
\end{proof}

\begin{proof}[Proof of Part~\ref{cor:ebs} of Corollary~\ref{corB}]
For every $p\gg0$, Proposition~\ref{prop:tau_p} gives $\mu(f_p)=\mu(f)<\infty$. Thus Corollary~\ref{cor:BS_F} and Lemma~\ref{lem:nu_E_BS} give $\fpt(f_p)\mu(f)\geq e^{\textup{BS}}(f)-1$. Since $\fpt(f_p)\leq\lct(f)$ for $p\gg0$, we obtain $\lct(f)\mu(f)+1\geq e^{\textup{BS}}(f)$.
\end{proof}

% REVIEW NOTE:
% The previous draft contained a coefficient-wise lifting lemma for plane branches,
% an alternative proof of Theorem~\ref{thm:semigr_fpt}, and a corollary derived from it.
% That argument has been removed because coefficient-wise lifting does not preserve
% intersection multiplicities in general; a correct replacement requires an
% equisingular-lifting/good-reduction theorem.

% \begin{remark}\label{rem:lift_semigroup}
% \textup{A coefficient-wise lift from characteristic $p$ to characteristic zero does not, in general, preserve intersection multiplicities, and therefore cannot by itself be used to preserve the semigroup of values. For example, over $\mathbb{F}_3[[x,y]]$, let
% \[
% f=y+2x,\qquad g=y^2+2x^2+x^3.
% \]
% Then $i(f,g)=3$, whereas for the coefficient-wise lifts
% \[
% \widetilde f=y+2x,\qquad \widetilde g=y^2+2x^2+x^3
% \]
% over $\mathbb{Q}$, one has $i(\widetilde f,\widetilde g)=2$. Thus any semigroup-preserving lifting argument requires additional equisingularity input.}
% \end{remark}

\begin{remark}
\begin{enumerate}
    \item \textup{The question of whether for every $f\in\mathbb{Z}[\underline{x}]$ there are infinitely many primes $p$ such that $\fpt(f_p)=\lct(f)$ is a fundamental conjecture in the study of $F$-pure thresholds; see Conjecture 3.6 in~\cite{mustata2005f}. It is known in some special cases (see, for example,~\cite{hernandez2016f,hernandez2011f,fedder1983F,smith1997vanishing,ein2007invariants}), but remains open in general.}
    \item \textup{Analogous results hold for a polynomial $f\in\mathbb{C}[\underline{x}]$ after spreading out its coefficients over a finitely generated $\mathbb{Z}$-algebra $A$. For every maximal ideal $\mathfrak n$ of $A$, the residue field $\frac{A}{\mathfrak n}$ is finite, hence isomorphic to $\mathbb{F}_{p^r}$ for some prime $p$ and some $r>0$. For more information, see Section 1 in Chapter 6 of~\cite{schwede2024singularities}. For isolated singularities, the determinacy theorem (Proposition~\ref{prop:mt_basic}) allows us to replace the relevant power series by contact-equivalent polynomials.}
\end{enumerate}
\end{remark}

\bibliographystyle{alpha}
\bibliography{bib}

\end{document}
